%! Author = Admin
%! Date = 17.04.2026
\section{Разработка архитектуры полносвязной нейронной сети}

Цель нейросетевого модуля состоит не в замене базового контроллера, а в адаптивном выборе параметрической скорости $V^{\ast}$ в зависимости от текущего состояния системы и локальной геометрии траектории. Тем самым нейросеть работает как интеллектуальная надстройка над аналитическим регулятором.

\subsection{Постановка задачи аппроксимации}

Требуется построить отображение
\begin{equation}
    \mathcal{F}: \mathbb{R}^{7} \rightarrow [V_{\min}, V_{\max}],
\end{equation}
которое по вектору признаков $f$ предсказывает скорость
\begin{equation}
    \hat{V}^{\ast} = \mathcal{F}(f),
\end{equation}
максимально близкую к значению $V_{\mathrm{opt}}$, найденному оракулом. Таким образом, решается задача регрессии
\begin{equation}
    \hat{V}^{\ast} \approx V_{\mathrm{opt}}.
\end{equation}

\subsection{Архитектура сети}

В проекте используется полносвязная нейронная сеть \texttt{SpeedMLP} со следующей структурой:
\begin{equation}
    7 \rightarrow 128 \rightarrow 128 \rightarrow 64 \rightarrow 1.
\end{equation}
После каждого скрытого линейного слоя применяется функция активации ReLU:
\begin{equation}
    \mathrm{ReLU}(x) = \max(0, x).
\end{equation}
Выходной слой формирует одно скалярное значение, которое затем ограничивается сверху через сигмоидную нелинейность:
\begin{equation}
    \hat{V}^{\ast} = V_{\max} \cdot \sigma(z),
\end{equation}
где
\begin{equation}
    \sigma(z) = \frac{1}{1 + e^{-z}}.
\end{equation}

Такое решение гарантирует, что сеть не предсказывает физически недопустимые скорости выше верхнего ограничения модели. На этапе инференса дополнительно выполняется явное отсечение результата к диапазону $[V_{\min},V_{\max}]$.

\subsection{Функция потерь и обучение}

Обучение выполняется как задача минимизации среднеквадратичной ошибки:
\begin{equation}
    \mathcal{L}_{\mathrm{MSE}} =
    \frac{1}{N}\sum_{i=1}^{N}\left(\hat{V}^{\ast}_i - V_{\mathrm{opt},i}\right)^2.
\end{equation}
В качестве метода оптимизации используется алгоритм Adam с параметрами по умолчанию:
\begin{equation}
    \eta = 10^{-3}, \qquad \text{batch size} = 64.
\end{equation}
Максимальное число эпох обучения равно $200$, однако фактически применяется ранняя остановка по валидационной ошибке:
\begin{equation}
    \text{stop, если } \mathcal{L}_{\mathrm{val}} \text{ не улучшается в течение } 20 \text{ эпох}.
\end{equation}

Использование ранней остановки снижает риск переобучения и позволяет сохранять лучшую версию весов сети по метрике \texttt{val\_loss}.

\subsection{Интеграция сети в систему управления}

После обучения сеть встраивается в контур управления через функцию \texttt{speed\_fn(state, s)}, возвращающую предсказанное значение $\hat{V}^{\ast}$. На каждом шаге симуляции происходит:
\begin{enumerate}
    \item вычисление признаков по текущему состоянию и параметру ближайшей точки;
    \item предсказание скорости нейросетью;
    \item ограничение скорости и ее производной;
    \item передача нового значения $V^{\ast}$ в базовый регулятор.
\end{enumerate}

С математической точки зрения это означает замену константы $V^{\ast}$ на адаптивный закон
\begin{equation}
    V^{\ast}(t) = \mathrm{clip}\bigl(\mathcal{F}(f(t)), V_{\min}, V_{\max}\bigr).
\end{equation}
Чтобы исключить резкие скачки, в коде дополнительно ограничивается темп изменения скорости:
\begin{equation}
    |V^{\ast}_{k+1} - V^{\ast}_{k}| \le a_{\max} \Delta t.
\end{equation}

\section{Визуализация результатов и выводы}

Система визуализации автоматически строит и сохраняет:
\begin{enumerate}
    \item трехмерную траекторию квадрокоптера и эталонной кривой;
    \item проекцию движения на плоскость $XY$;
    \item графики ошибок $s_{\mathrm{arc}} - V^{\ast} t$, $e_1$, $e_2$;
    \item график ошибки по курсу $\delta_{\phi}$;
    \item график нормы линейной скорости $\|v\|$;
    \item графики углов $\phi$, $\theta$, $\psi$;
    \item графики обучения нейросети по эпохам.
\end{enumerate}

Такая визуализация решает две задачи. Во-первых, она позволяет качественно оценить сходимость на траектории разных типов. Во-вторых, она делает возможным прямое сравнение двух режимов: с постоянной скоростью $V^{\ast}=\mathrm{const}$ и с адаптивной скоростью, формируемой нейросетью.

По смыслу работы ожидаемый эффект от нейросетевого модуля состоит в следующем:
\begin{enumerate}
    \item на участках малой кривизны сеть может повышать скорость без потери устойчивости;
    \item на сложных участках с большой кривизной или неблагоприятной ориентацией скорость должна уменьшаться;
    \item базовый аналитический контроллер при этом сохраняет устойчивость и отвечает за само слежение, а нейросеть только улучшает режим движения.
\end{enumerate}

Следовательно, интеллектуальный модуль не заменяет математическую модель и синтезированный регулятор, а расширяет их, обеспечивая адаптацию к текущей геометрии задачи.

На рисунках~\ref{fig:train_loss}--\ref{fig:elliptic_errors} представлены результаты обучения нейросети и работы системы на эллиптической спирали.

\begin{figure}[htbp]
    \centering
    \begin{subfigure}[b]{0.48\textwidth}
        \centering
        \includegraphics[width=\textwidth]{images/training_loss}
        \caption{Кривые обучения SpeedMLP}
        \label{fig:train_loss}
    \end{subfigure}
    \hfill
    \begin{subfigure}[b]{0.48\textwidth}
        \centering
        \includegraphics[width=\textwidth]{images/prediction_quality}
        \caption{Качество предсказания $\hat{V}^{\ast}$ vs $V_{\mathrm{opt}}$}
        \label{fig:pred_quality}
    \end{subfigure}
    \caption{Результаты обучения нейросетевого оптимизатора скорости}
    \label{fig:nn_training}
\end{figure}

\begin{figure}[htbp]
    \centering
    \begin{subfigure}[b]{0.48\textwidth}
        \centering
        \includegraphics[width=\textwidth]{images/elliptic_traj_3d}
        \caption{Траектория, эллиптическая спираль}
        \label{fig:elliptic_3d}
    \end{subfigure}
    \hfill
    \begin{subfigure}[b]{0.48\textwidth}
        \centering
        \includegraphics[width=\textwidth]{images/elliptic_errors}
        \caption{Ошибки слежения}
        \label{fig:elliptic_errors}
    \end{subfigure}
    \caption{Результаты симуляции на эллиптической спирали с нейросетевым оптимизатором $V^{\ast}$}
    \label{fig:elliptic_results}
\end{figure}
